\documentclass[11pt]{article}
\usepackage{amsmath,amssymb,amsthm,geometry,hyperref}
\geometry{margin=1in}
\newcommand{\C}{\mathbb C}
\newcommand{\Span}{\operatorname{span}}
\newcommand{\diag}{\operatorname{diag}}
\title{Complete native differential at every finite multiple-root unfolding}
\author{Independent programme derivation}
\date{20 September 2026}
\begin{document}\maketitle
The received exceptional atlas \cite{Atlas}, EA29--36, specifies
constant-term unfoldings and claims their original-metric singular
asymptotes. This note verifies those claims, derives the complete
forward and inverse compound limits with their complex phases,
evaluates every multiplicity-dependent constant, and determines
the original singular flags. The polynomial is the retained
four-dimensional ES construction \cite{ES}; its existing complete
inverse is GF7. Every metric below is the fixed original Gamma
metric transported by the existing \(\Psi\), with no change to
the period or the conductor as the unfolding parameter varies.

\section{The entire original differential and its inverse}
Use the original source coordinates \((a,y,z,w)\), and output
order \((A,B,C,D)\). The polynomial is
\[
\begin{aligned}
P_0&=a^3z+2a^2y-ia,\\
P_2&=-a^3y^2z-2ia^2yz+a^2w-2a^2y^3-10iay^2+3az+y,\\
P_3&=2a^3y^3z+6ia^2y^2z+2a^2wy+4a^2y^4
       +2iaw-4iay^3-2ayz+2iz+7y^2,\\
P_4&=2a^3y^4z+8ia^2y^3z+a^2wy^2+4a^2y^5
       +2iawy+7iay^4-10ay^2z-4iyz-w-3y^3.
\end{aligned}\tag{FD1}
\]
For \(t\ne0\), put
\[
\begin{aligned}
q(t,r,A,B)=\bigl(&t^{-1},-r-it,At^3+2rt+3it^2,\\
 &7ir^2t+(B-17r+Ar^2)t^2-13it^3-2At^4\bigr)^T,\\
\pi(t,r,A,B)=\bigl(&A,B,t^2-4Ar^3-3r^2-2Br,\,
                      3Ar^4+2r^3+Br^2-rt^2\bigr)^T.
\end{aligned}\tag{FD2}
\]
Substitution into FD1 proves \(Pq=\pi\). This is the complete
finite simple-root inverse, with \(t=1/a\), rather than a
replacement source coordinate at \(t=0\). With
\[
 h_u(s)=As^4+s^3+Bs^2+Cs+D,\quad
 \ell_r=(r^4,r^2,r,1),\quad \ell'_r=(4r^3,2r,1,0),\quad
 \kappa=12Ar^2+6r+2B,
 \tag{FD3}
\]
its equations are \(h_u(r)=0\), \(t^2=h_u'(r)\).
Differentiating these two equations in all four output directions
gives exactly
\[
 dr=-\ell_r\,du/t^2,\qquad
 dt=\ell'_r\,du/(2t)-\kappa\ell_r\,du/(2t^3).
 \tag{FD4}
\]
Here and throughout, \(\ell_r\) is the complex-linear row
functional with the displayed coefficients, not their conjugates.
The differential of the original coordinate vector in FD2 is
\[
\begin{aligned}
 da&=-dt/t^2,&dy&=-dr-i\,dt,\\
 dz&=t^3dA+2t\,dr+(3At^2+2r+6it)dt,\\
 dw&=(r^2t^2-2t^4)dA+t^2dB
       +[14irt+(-17+2Ar)t^2]dr\\
   &\hspace{13mm}
       +[7ir^2+2(B-17r+Ar^2)t-39it^2-8At^3]dt.
\end{aligned}\tag{FD5}
\]
FD4 inserted into FD5 is the full \(4\)-by-\(4\) matrix
\(J^{-1}=DP(q)^{-1}\), in the original coordinate orders.
In particular it includes every \(dA,dB,dC,dD\) term.
An equivalent explicit outer-product formula is
\[
 J^{-1}=q_t\left(\frac{\ell'_r}{2t}
                    -\frac{\kappa\ell_r}{2t^3}\right)
          -\frac{q_r\ell_r}{t^2}+q_Ae_A^*+q_Be_B^*,
 \tag{FD6}
\]
where every column \(q_t,q_r,q_A,q_B\) is specified by FD5.
Thus FD6 has no unevaluated matrix inverse.

For completeness, the entire forward Jacobian also has a compact
column expression retaining all lower terms. Put
\[
\begin{aligned}
v(r)&=(1,-r^2,-2r^3,2r^4)^T,&
w(r)&=(0,1,-2r,r^2)^T,\\
R_0&=(4i,3ir^2,-22ir^3,15ir^4)^T,\\
R_1&=(3A,-Ar^2+2B-16r,
                  -10Ar^3-4Br+32r^2,8Ar^4+2Br^2-16r^3)^T,\\
R_2&=(0,-2iAr-14i,-8iAr^2-2iB+22ir,
                                10iAr^3+2iBr-8ir^2)^T,\\
R_3&=(0,-2A,4Ar-2,-2Ar^2+2r)^T,\\
S_0&=(0,2Ar-13,8Ar^2+2B+32r,
                              -10Ar^3-2Br-19r^2)^T.
\end{aligned}
\]
Writing \(e_4=(0,0,0,1)^T\), its four columns are exactly
\[
\begin{aligned}
J_{\bullet a}&=\frac{2r}{t}v+R_0+tR_1+t^2R_2+t^3R_3-it^4e_4,\\
J_{\bullet y}&=\frac2{t^2}v+\frac{14ir}{t}w+S_0+t^2e_4,\\
J_{\bullet z}&=\frac1{t^3}v+\frac2t w,\qquad
J_{\bullet w}=\frac1{t^2}w .
\end{aligned}\tag{FD7}
\]
These columns follow by differentiating FD1 before substituting
FD2. The independent certificate verifies both products of FD7
with FD4--5 to be the identity over
\(\mathbb Q(i)(t,r,A,B)\).

The two intermediate Jacobians have determinants
\[
 \det Dq=t^3,\qquad \det D\pi=-2t^3,\qquad \det J=-2.
 \tag{FD8}
\]
For the first formula, after the first two rows the diagonal
entries of the remaining triangular minor are \(t^3,t^2\);
the first two pivots are \(-t^{-2},-1\). For the second,
move the two coefficient columns \(A,B\) to the front, an
even permutation. The residual \(t,r\) block is
\(\left(\begin{smallmatrix}2t&-\kappa\\-2rt&r\kappa-t^2\end{smallmatrix}\right)\),
whose determinant is \(-2t^3\). The chain rule proves the
last equality for every \(t\ne0\). These parameter values
cover the original \(a\ne0\) chart by \(t=1/a\),
\(r=-y-i/a\), \(A=P_0\), \(B=P_2\). Polynomial continuation
therefore gives \(\det DP=-2\) everywhere in the original
affine source. No finite critical point is created here.

\section{The exact complex Puiseux branches}
Fix an actual target polynomial \(h_0(s)=As^4+s^3+Bs^2+Cs+D\)
with a finite root \(r_0\) of exact multiplicity \(m=2,3,4\).
For \(x=s-r_0\), its original coefficients give exactly
\[
 h_0(r_0+x)=c_2x^2+c_3x^3+c_4x^4,\qquad
 c_2=6Ar_0^2+3r_0+B,\quad c_3=4Ar_0+1,\quad c_4=A.
 \tag{FD9}
\]
The first nonzero coefficient is \(c_m\).
For \(m=4\), in particular, \(c_3=0\) forces
\(A=-1/(4r_0)\) and \(r_0\ne0\); these are consequences
of the retained cubic coefficient one, not free normalizations.
The specified unfolding is exactly
\[
 h_\varepsilon(s)=h_0(s)-\varepsilon,\qquad
 u_\varepsilon=u_0-\varepsilon e_D.
 \tag{FD10}
\]
Keep the actual complex \(c_m\ne0\). The analytic map
\[
 f(x)=x\left(\frac{h_0(r_0+x)}{c_mx^m}\right)^{1/m}
 \tag{FD11}
\]
uses the unique local root with value one at \(x=0\).
Its derivative is one there, so it has a local analytic inverse
\(g(w)=w+O(w^2)\). This follows, for example, by the complex
inverse-function theorem; its nonzero derivative is explicit.
Let \(\zeta_j=e^{2\pi ij/m}\), \(0\le j<m\), and use the
Puiseux parameter \(u\) with \(\varepsilon=c_mu^{2m}\).
Then all \(m\) local roots and both original factor signs are
given by the exact analytic branches
\[
\begin{aligned}
 r_j(u)&=r_0+g(\zeta_j u^2),\\
 t_{j,\pm}(u)&=\theta_{j,\pm}u^{m-1}a_j(u),\qquad
 \theta_{j,\pm}^2=mc_m\zeta_j^{m-1},\qquad
 \theta_{j,-}=-\theta_{j,+},
\end{aligned}\tag{FD12}
\]
where \(a_j(u)\) is the analytic square root with value one of
\[
 \frac{h_0'(r_j(u))}
                  {mc_m\zeta_j^{m-1}u^{2m-2}}.
\]
The quotient has analytic extension with value one. In particular
\[
\begin{aligned}
 r_j-r_0&=\zeta_ju^2+O(u^4),\\
 t_{j,\pm}&=\theta_{j,\pm}u^{m-1}(1+O(u^2)),\\
 \kappa_j=h_0''(r_j)&=K_j u^{2m-4}(1+O(u^2)),\qquad
 K_j=m(m-1)c_m\zeta_j^{m-2},\\
 t_{j,\pm}^2/\kappa_j&=\frac{\zeta_j}{m-1}u^2(1+O(u^2)).
\end{aligned}\tag{FD13}
\]
These formulas hold as complex \(u\to0\), on either sign
branch, with every phase in \(\theta_{j,\pm},K_j,\zeta_j\)
retained. For sufficiently small nonzero \(u\), both \(t\)
and \(\kappa\) are nonzero. All resulting
\(q(t_{j,\pm}(u),r_j(u),A,B)\) are actual original source
points, by FD2 and FD10--12. No claim about a different
target path is used.

\section{All phased forward and inverse exterior limits}
Use increasing coordinate subsets of \(\{1,2,3,4\}\), in
the order \(a,y,z,w\) in the domain and \(A,B,C,D\) in the
codomain. A superscript star on a coordinate wedge \(e_I^*\)
in the following matrices means its coordinate covector;
it is not a metric adjoint or a conjugation of the root.
Let \(v_0=v(r_0)\), \(w_0=w(r_0)\),
\(\ell_0=(r_0^4,r_0^2,r_0,1)\), and
\(\ell'_0=(4r_0^3,2r_0,1,0)\). Define
\[
\begin{aligned}
 b_3&=(1,-r_0,r_0^2,-r_0^4)^T
       &&\text{in the wedge order }123,124,134,234,\\
 \eta_3&=(r_0^4,r_0^3,-r_0^2/2,-1/2)
       &&\text{in that same wedge order}.
\end{aligned}
\]
Along each branch FD12 the complete matrix limits are
\[
\boxed{\begin{array}{c|c|c}
 r&\text{forward limit}&\text{inverse limit}\\\hline
 1&t^3J\longrightarrow v_0e_3^*
       &(t^5/\kappa)J^{-1}\longrightarrow\tfrac12e_1\ell_0\\
 2&t^5\wedge^2J\longrightarrow(v_0\wedge w_0)e_{34}^*
       &t^5\wedge^2J^{-1}\longrightarrow
                    -\tfrac12e_{12}(\ell'_0\wedge\ell_0)\\
 3&(t^5/\kappa)\wedge^3J\longrightarrow b_3e_{234}^*
       &t^3\wedge^3J^{-1}\longrightarrow e_{124}\eta_3\\
 4&\wedge^4J=-2&\wedge^4J^{-1}=-\tfrac12
\end{array}.}\tag{FD14}
\]
Every entry and sign can also be read as the following
explicit wedge coefficient rows:
\[
\begin{aligned}
 v_0\wedge w_0&=(1,-2r_0,r_0^2,4r_0^3,-3r_0^4,2r_0^5)^T,\\
 \ell'_0\wedge\ell_0&=(2r_0^5,3r_0^4,4r_0^3,r_0^2,2r_0,1).
\end{aligned}\tag{FD15}
\]
For \(r=1\) on the forward side, FD7 gives the limit
entry by entry. For \(r=2\), the only order \(t^{-5}\)
minor uses columns \(z,w\). The potential order \(t^{-5}\)
term from columns \(y,z\) vanishes because both leading
columns are multiples of \(v(r)\); all remaining terms
have smaller pole. This proves the full second-compound
limit, not only one of its entries.

For the inverse first-compound limit, its first row in
FD4--5 is exactly
\[
 \frac{\kappa}{2t^5}\ell_r-\frac1{2t^3}\ell'_r.
 \tag{FD16}
\]
Multiplication by \(t^5/\kappa\), followed by FD13,
gives \(\ell_0/2\). The other rows tend to zero under
this factor: their terms contain positive powers of \(t\),
or such powers times \(t^2/\kappa\), after the multiplication.
For example the \(y\)-row becomes
\(t^3\ell_r/\kappa-it^4\ell'_r/(2\kappa)
+i t^2\ell_r/2\), which tends to zero. The \(z,w\)
rows follow directly from their bounded coefficient
polynomials in FD5 and the same relation
\(t^2/\kappa=O(u^2)\). This proves the inverse limit
also when \(\kappa\to0\), as happens for \(m=3,4\).

All remaining limits follow from the exact complementary
minor identity, with its original determinant:
\[
 \det(J^{-1})_{I,J}
 =(-1)^{\sum I+\sum J}\frac{\det J_{J^c,I^c}}{\det J}
 =-\tfrac12(-1)^{\sum I+\sum J}\det J_{J^c,I^c}.
 \tag{FD17}
\]
This identity follows from the adjugate formula, or by
expanding the determinant on complementary coordinate
subsets. For the third forward compound use the same
identity for \(J^{-1}\) and the first inverse limit.
It gives the column \((1,-r_0,r_0^2,-r_0^4)^T\)
at column \(234\). For the second inverse compound
it gives only row \(12\), with the row in FD15 multiplied
by \(-1/2\). For the third inverse compound it gives
only row \(124\), equal to \(\eta_3\). Thus the signs
and coordinate complements in FD14 are fully determined.

One can state these limits directly in the complex Puiseux
parameter without hiding any phase. For a fixed \(j,\pm\)
write \(\theta=\theta_{j,\pm}\), \(K=K_j\). The forward
scalings and coefficients are
\[
\begin{aligned}
 u^{3(m-1)}J&\longrightarrow\theta^{-3}v_0e_3^*,\\
 u^{5(m-1)}\wedge^2J&\longrightarrow
                     \theta^{-5}(v_0\wedge w_0)e_{34}^*,\\
 u^{3m-1}\wedge^3J&\longrightarrow K\theta^{-5}b_3e_{234}^*,
\end{aligned}\tag{FD18}
\]
and the inverse limits are respectively
\[
\begin{aligned}
 u^{3m-1}J^{-1}&\longrightarrow
                         \frac{K}{2\theta^5}e_1\ell_0,\\
 u^{5(m-1)}\wedge^2J^{-1}&\longrightarrow
                         -\frac1{2\theta^5}e_{12}
                                      (\ell'_0\wedge\ell_0),\\
 u^{3(m-1)}\wedge^3J^{-1}&\longrightarrow
                         \theta^{-3}e_{124}\eta_3.
\end{aligned}\tag{FD19}
\]
These give the phase for every one of the \(2m\) signed
branches, including the effect of \(t\mapsto-t\).

\section{The original metric and all constants}
Fix the original receiving map
\[
 \Psi=B_v(V_\rho^{\mathsf T})^{-1}K_{\rm ES}^{-1},\qquad
 (B_v)_{kj}=\begin{cases}
 \binom{v+j}{k}\mu_{v+j-k},&k\le j,\\0,&k>j,
 \end{cases}\quad 0\le k,j\le3,
 \tag{FD20}
\]
with its original fixed conductor order, moments,
quartet Vandermonde and collision matrix. Its determinant
is nonzero by the already proved FC27 and the original
\(\mu_v\ne0\). The norm on the target polynomial space
is unchanged:
\[
 \|p\|_\Gamma^2=\int_\R|p(c_*+iy)|^2dm_s(y),\qquad
 dm_s(y)=\frac{(2\pi)^{s/2}2^{s/2}}{4\pi\Gamma(s/2)}
                |\Gamma(s/4+iy/2)|^2dy .
 \tag{FD21}
\]
Let \(G=\Psi^*\mathcal H_\Gamma\Psi\) be its full positive
pullback. In the original monomial frame every entry is
explicit:
\[
 (\mathcal H_\Gamma)_{jk}
 =\sum_{p=0}^j\sum_{q=0}^k
 \binom jp\binom kq c_*^{j+k-p-q}(-i)^pi^q m_{p+q},
 \tag{FD22}
\]
where \(b=s/2\), \(M=(2\pi)^{s/2}\),
\[
 m_0=M,\quad m_2=Mb,\quad m_4=M(3b^2+2b),\quad
 m_6=M(15b^3+30b^2+16b),\qquad m_1=m_3=m_5=0.
\]
These are the original WCF Gamma moments \cite{WCF}.
All period coefficients, phases, center, mass and cross
terms remain in FD20--22. The fixed \(\Psi\) conjugates
the polynomial \(P\) to \(P_W=\Psi P\Psi^{-1}\); therefore
the derivative \(J\) uses the same \(G\) on its domain
and codomain. No moving-period limit is taken.

Define the three strictly positive original constants
\[
\begin{aligned}
 C_1&=\sqrt{v_0^*Gv_0}\sqrt{(G^{-1})_{33}},\\
 C_2&=\sqrt{\det([v_0\ w_0]^*G[v_0\ w_0])}
                       \sqrt{\det(G^{-1})_{34,34}},\\
 C_3&=\sqrt{G_{11}}\sqrt{\ell_0G^{-1}\ell_0^*}.
\end{aligned}\tag{FD23}
\]
Positivity follows from \(v_{0,1}=1\), independence
of \(v_0,w_0\), \(\ell_{0,4}=1\), and positive definiteness
of \(G\). The norm of a rank-one exterior map
\(x\lambda\) is \(\|x\|\|\lambda\|\): Cauchy--Schwarz
gives the upper bound and the Riesz representative of
\(\lambda\) attains it. The Gram determinant formula
for exterior norms therefore gives \(C_1,C_2\) for the
first two forward limits and \(C_3/2\) for the first
inverse limit in FD14.

For an invertible endomorphism of this fixed four-dimensional
Hermitian space, products of singular values prove the exact
identity
\[
 \|\wedge^rJ^{-1}\|_G
            =\frac{\|\wedge^{4-r}J\|_G}{|\det J|}
            =\tfrac12\|\wedge^{4-r}J\|_G .
 \tag{FD24}
\]
Indeed the inverse singular values are the reciprocals
in reverse order; canceling their product proves the
formula. Consequently the complete exterior asymptotes
in the original norm are
\[
\begin{array}{c|c|c}
 r&\|\wedge^rJ\|_G&\|\wedge^rJ^{-1}\|_G\\\hline
 1&\sim C_1|t|^{-3}&\sim C_3|\kappa|/(2|t|^5)\\
 2&\sim C_2|t|^{-5}&\sim C_2/(2|t|^5)\\
 3&\sim C_3|\kappa|/|t|^5&\sim C_1/(2|t|^3)\\
 4&2&1/2
\end{array}.\tag{FD25}
\]
The limits are norm limits of the full finite matrices,
not estimates from sampled numerical singular values.

Successive ratios of forward exterior norms give every
singular value in decreasing order:
\[
\boxed{\quad
 \sigma_1\sim C_1|t|^{-3},\qquad
 \sigma_2\sim\frac{C_2}{C_1}|t|^{-2},\qquad
 \sigma_3\sim\frac{C_3}{C_2}|\kappa|,\qquad
 \sigma_4\sim\frac{2|t|^5}{C_3|\kappa|}.
 \quad}\tag{FD26}
\]
For completeness \(\|\wedge^rJ\|\) is the product of the
largest \(r\) singular values: diagonalizing \(J^*J\) in
the fixed metric and taking wedges gives all \(r\)-fold
products. This proves the ratio step and verifies EA36.

The exact coefficient \(c_m\) gives
\[
 |t|\sim\sqrt m\,|c_m|^{1/(2m)}
                       |\varepsilon|^{(m-1)/(2m)},\qquad
 |\kappa|\sim m(m-1)|c_m|^{2/m}
                       |\varepsilon|^{(m-2)/m}.
 \tag{FD27}
\]
Substituting, with no amplitude set to one, yields
\[
\boxed{\begin{aligned}
 \sigma_1&\sim\frac{C_1}{m^{3/2}|c_m|^{3/(2m)}}
                  |\varepsilon|^{-3(m-1)/(2m)},\\
 \sigma_2&\sim\frac{C_2}{mC_1|c_m|^{1/m}}
                  |\varepsilon|^{-(m-1)/m},\\
 \sigma_3&\sim\frac{m(m-1)C_3|c_m|^{2/m}}{C_2}
                  |\varepsilon|^{(m-2)/m},\\
 \sigma_4&\sim\frac{2m^{3/2}|c_m|^{1/(2m)}}{(m-1)C_3}
                  |\varepsilon|^{(3m-1)/(2m)}.
\end{aligned}}\tag{FD28}
\]
All four constants multiply to exactly \(2\), agreeing
with FD8, and every phase has already been recorded in
the stronger complex matrix limits FD18--19. The powers are
\[
\begin{array}{c|rrrr}
m&\sigma_1&\sigma_2&\sigma_3&\sigma_4\\\hline
2&-3/4&-1/2&0&5/4\\
3&-1&-2/3&1/3&4/3\\
4&-9/8&-3/4&1/2&11/8
\end{array}.
\]
They belong to the specified constant-term path FD10.

\section{Original singular flags and exact fixed-plane volumes}
Every adjacent singular-value ratio tends to infinity
by FD28. Thus the leading singular flags are unique for
sufficiently small nonzero parameter. The leading right
one-, two-, and three-planes converge, in the original
metric, to
\[
 G^{-1}\Span(e_3),\qquad
 G^{-1}\Span(e_3,e_4),\qquad
 G^{-1}\Span(e_2,e_3,e_4).
 \tag{FD29}
\]
The corresponding leading left planes converge to
\[
 \Span(v_0),\qquad \Span(v_0,w_0),\qquad \ker\ell_0.
 \tag{FD30}
\]
For the proof, divide each full forward compound by its
norm and apply FD14. Its limit is a unit-norm rank-one
map, so its products with its adjoint tend to rank-one
orthogonal projections. The top eigenspace projections
converge: any subsequential limit of a unit top eigenvector
satisfies the limiting eigenvector equation and lies in
the unique limiting range. The right Riesz representative
of the coordinate wedge functional \(e_J^*\) is
\(\bigwedge_{j\in J}G^{-1}e_j\). The left third compound
vector \(b_3\) is decomposable and represents \(\ker\ell_0\),
as follows by taking its wedge coefficients or the
complementary-minor formula FD17. This proves both flags.
In particular the smallest right singular line tends to
\(\C e_1\), and the smallest left line to
\(\C G^{-1}\ell_0^*\), by taking their original-metric
orthogonal complements. Their actual polynomial-space
images are obtained by the fixed \(\Psi\).

The exact volume distortion on the fixed coordinate
input flags \(\Span(e_3)\), \(\Span(e_3,e_4)\), and
\(\Span(e_2,e_3,e_4)\) is respectively
\[
\begin{aligned}
 &\sim |t|^{-6}\frac{v_0^*Gv_0}{G_{33}},\\
 &\sim |t|^{-10}
              \frac{\det([v_0\ w_0]^*G[v_0\ w_0])}{\det G_{34,34}},\\
 &\sim |\kappa|^2|t|^{-10}
              \frac{\ell_0G^{-1}\ell_0^*}{(G^{-1})_{11}}.
\end{aligned}\tag{FD31}
\]
Each is a squared output Gram determinant divided by
the squared input volume in the original metric. For
the first two apply the compound limit to the displayed
coordinate wedge. For the third, the complementary
Gram identity gives
\(\|b_3\|_{\wedge^3G}^2
=\det G\,\ell_0G^{-1}\ell_0^*\) and
\(\det G_{234,234}=\det G(G^{-1})_{11}\).
This proves the last constant including its cross terms.
In degree four the squared-volume ratio is exactly four
at every finite point.

\section{Escape and the exact conductor relation}
Formula FD2 and the actual branches FD12 imply that
\(a=1/t\) is the sole divergent coordinate; the other
three tend to \((-r_0,0,0)\). Thus
\[
 |\varepsilon|^{(m-1)/(2m)}\|\Psi q\|_\Gamma
 \longrightarrow\frac{\sqrt{G_{11}}}
                         {\sqrt m\,|c_m|^{1/(2m)}}>0.
 \tag{FD32}
\]
The original factor is \(L=a(U-rV)\), so its actual
coefficient \(a\) diverges as well. This is intrinsic
factor-scale escape at the selected finite repeated root;
adding a missing \(a=0\) chart cannot supply a finite
limit of that coordinate. The image displacement is
exactly \(-\varepsilon\Psi e_4\). Consequently, if
\(\alpha_m=(m-1)/(2m)\),
\[
 \|\Psi(u_\varepsilon-u_0)\|_\Gamma^{\alpha_m}
 \|\Psi q\|_\Gamma
 \longrightarrow
 \frac{G_{44}^{\alpha_m/2}\sqrt{G_{11}}}
                         {\sqrt m\,|c_m|^{1/(2m)}}.
 \tag{FD33}
\]

The fixed original source map \(\Phi\) and conductor
restriction \(L=T_A|_U\) obey \(L\Phi=\Psi\).
For \(P_U=\Phi P\Phi^{-1}\), \(P_W=\Psi P\Psi^{-1}\),
differentiation at the matching points gives exactly
\[
 L\,DP_U(\Phi q)=DP_W(\Psi q)L.
 \tag{FD34}
\]
Both maps use their original respective norms; the
source-side formulas follow by replacing \(G\) by
the actual fixed \(G_U=\Phi^*\mathcal H_U\Phi\), with
no change to the differential matrices or phases.
In particular, exterior inverse transport is the exact
identity
\[
 \wedge^r DP_W^{-1}
   =(\wedge^rL)(\wedge^rDP_U^{-1})(\wedge^rL^{-1}).
 \tag{FD35}
\]
This is the original four-plane inverse from FC24, not
a variable conductor matrix or a claim that its leading
moment vanished. The finite-root degenerations above
are completely located in the specified nonlinear
derivative along its actual escaping branches.

\section{Exact verification and source use}
The independent script
\texttt{verify\_finite\_root\_native\_differential.py}
starts from FD1, differentiates it before substitution,
checks \(Pq=\pi\), both intermediate determinants,
\(JDq=D\pi\), both inverse products for FD4--5,
and every forward and inverse compound over
\(\mathbb Q(i)(t,r,A,B)\). Its full matrices and checks
are recorded in
\texttt{finite\_root\_native\_differential\_exact.json}.
No sampled period, artificial Gram or numerical
singular-value diagnostic is used in any proof above.
The Puiseux phases and the passage to the original metric
are proved explicitly in FD9--33.

\begin{thebibliography}{9}
\bibitem{Atlas} User-supplied continuation,
\emph{Exceptional-map atlas: the finite signed-root completion},
\texttt{EXCEPTIONAL\_ATLAS.md}, 20 September 2026,
EA29--45, especially EA33--36; no author name was supplied.
Exact received-source SHA-256:
\texttt{fbe6de554732ccd911fe36e75b7d566bd5dbf28719592c8d6a5f56e177c9fafd}.
The present note independently verifies the claimed asymptotes
and adds the full phased inverse compounds, explicit \(c_m\)
constants and original singular flags.
\bibitem{ES} The Clankers, \emph{Erd\H{o}s--Straus project reader},
canonical 488-page edition, source version 69, theorem
\texttt{explicit-four-time-Jacobian-collision}, original
source lines 19484--19588; SHA-256
\texttt{678b3f70ee7d2a6d00f6c588d7778e390118ab4743a5759612dbba6430984c78}.
The original three-dimensional example is credited to
Levent Alp\"oge and Fable in Terence Tao's
\emph{A digestion of the Jacobian conjecture counterexample},
21 July 2026; the explicit four-dimensional marked ES
construction is the polynomial retained here.
\bibitem{FC} Split-Zero programme,
\texttt{FABLE\_TO\_ORIGINAL\_CONDUCTOR.tex},
FC24--35 and GF1--11; the exact fixed source and target
maps, quotient norms and original complete inverse.
\bibitem{WCF} Split-Zero programme,
\texttt{WEIGHTED\_CONDUCTOR\_FORWARD.tex}, WCF1--6,
crediting T. H. Koornwinder, R. Wong, R. Koekoek,
R. F. Swarttouw and W. P. Reinhardt, NIST DLMF
Chapter 18, 18.19.8--9, 18.22.8 and 18.23.7.
The original mass, center and phases are retained in FD21--22.
\end{thebibliography}
\end{document}
